\section{Dynamics, Endpoints, and Trend Sensitivity}\label{app:dynamics}

\subsection{Quarterly system}

The corrected dynamic model interacts every nonreference exposure quintile with calendar quarter.  The omitted period is 2022Q4, which contains October and November because December 2022 is the transition month and is excluded.  The resulting Q5--Q1 path has 23 preperiod and 15 postperiod coefficients.  The pooled and SOC2-by-calendar-month versions are estimated on identical reconstructed support, and the archive stores the complete Q2--Q5 coefficient vector, full covariance matrix, influence matrix, and simultaneous critical values.

Relative to omitted 2022Q4, the joint null that all 23 Q5 preperiod
coefficients equal zero is rejected in the pooled model ($p=.0171$) and the
family-conditioned model ($p=.0053$).  Table~\ref{tab:app_pretrend_windows}
also reports the reference-invariant null that the coefficients are equal to
one another within each fixed window.  That null rejects over the full period
for the pooled model ($p=.0121$) but not for the family-conditioned model
($p=.0641$); it does not reject at five percent in 2017Q1--2019Q4,
2021Q1--2022Q3, or the full window excluding 2020Q2--2020Q4 for either
structure.  Thus the full pooled rejection no longer clears five percent when
2020Q2--2020Q4 are excluded, whereas the family reference-relative rejection
persists even though equality within each fixed window is not rejected.  In
leave-one-quarter diagnostics,
2019Q2 produces the largest reduction in the pooled reference-relative Wald
statistic and 2018Q4 the largest reduction in the conditioned statistic; these
nonadditive deletions do not identify unique causal drivers.

\begin{table}[!htbp]
\centering
\caption{Preperiod joint tests across fixed calendar windows}\label{tab:app_pretrend_windows}
\resizebox{\textwidth}{!}{%
\begin{tabular}{llrr}
\toprule
Structure & Window & Equal to 2022Q4 & Equal within window \\
\midrule
Pooled & 2017Q1--2022Q3 & 0.0171 & 0.0121 \\
Pooled & 2017Q1--2019Q4 & 0.1365 & 0.0925 \\
Pooled & 2021Q1--2022Q3 & 0.4285 & 0.3281 \\
Pooled & Full, excluding 2020Q2--Q4 & 0.0772 & 0.0676 \\
\addlinespace
Family by month & 2017Q1--2022Q3 & 0.0053 & 0.0641 \\
Family by month & 2017Q1--2019Q4 & 0.0126 & 0.0744 \\
Family by month & 2021Q1--2022Q3 & 0.0116 & 0.3264 \\
Family by month & Full, excluding 2020Q2--Q4 & 0.0079 & 0.1048 \\
\bottomrule
\end{tabular}}
\tabnote{Entries are common-multiplier $p$-values.  The third column tests every selected preperiod coefficient against the omitted 2022Q4 coefficient, which represents October--November 2022.  The fourth tests equality among coefficients inside the selected window and is invariant to the omitted reference.  The two nulls therefore answer different questions.  Windows were declared before this result audit and are not selected for favorable $p$-values.}
\end{table}


As a precision diagnostic, an OLS linear functional through all preperiod
coefficients is $-0.0078$ log point per year in the pooled model and $0.0041$
after family conditioning.  Its occupation-cluster standard errors are 0.0119
and 0.0292, implying normal-theory $\mathrm{MDE}_{80}$ values of 0.0332 and
0.0819 log point per year.  An insignificant linear slope cannot rule out
economically relevant smooth drift and is not a substitute for the joint
tests.  Simultaneous bands exclude zero only in late pooled quarters and never
in the family-conditioned path.  These results do not establish which
preperiod force produces the differences and do not validate parallel
evolution for a causal interpretation of the January 2023 split.

Table~\ref{tab:app_static_dynamic_reconciliation} distinguishes three targets.
The postperiod functional $P$ averages the 15 post-quarter coefficients with
predeclared calendar weights relative to omitted 2022Q4.  Its pooled value is
$-0.1199$, with interval $[-0.2637,0.0239]$; its family-conditioned value is
$-0.2074$, with interval $[-0.4252,0.0103]$.  The reference-invariant dynamic
functional $D$ subtracts the corresponding weighted preperiod mean.  It is
$-0.131372$ pooled and $-0.021689$ after family conditioning, compared with
static coefficients $S$ of $-0.132109$ and $-0.021675$.  Covariance-preserving
tests give $p=.454$ and $p=.995$ for $D-S$.  Hence the static and dynamic
post-minus-pre targets are numerically reconciled.  Within the family model,
$P-S=-0.1858$ has $p=.074$, so the design does not detect that difference at
five percent.  Across conditioning structures, however, the change in $P-S$ is
$-0.1980$ ($p=.021$), whereas the change in $D-S$ is $-0.0008$ ($p=.693$).
Thus the short-reference functional responds to conditioning differently from
the reference-invariant comparison.  The $D-S$ result is a nondetection, not an
equivalence test, and $D$ remains a linear functional rather than the nonlinear
static grouped-binomial coefficient.

\begin{table}[!htbp]
\centering
\caption{Static and dynamic target reconciliation}\label{tab:app_static_dynamic_reconciliation}
\resizebox{\textwidth}{!}{%
\begin{tabular}{lrrrrr}
\toprule
Specification & Static $S$ & Post/reference $P$ & Post/pre $D$ & $D-S$ & $p(D-S=0)$ \\
\midrule
Pooled & $-0.132109$ & $-0.119889$ & $-0.131372$ & $\phantom{-}0.000737$ & 0.454 \\
Family by calendar month & $-0.021675$ & $-0.207434$ & $-0.021689$ & $-0.000014$ & 0.995 \\
Family minus pooled & $\phantom{-}0.110434$ & $-0.087545$ & $\phantom{-}0.109683$ & $-0.000752$ & 0.693 \\
\bottomrule
\end{tabular}}
\tabnote{$S$ is the static grouped-binomial Q5-by-post coefficient. $P$ is the observed-post-month-weighted average of quarterly Q5 coefficients relative to the two observed months in omitted 2022Q4. $D$ subtracts the corresponding observed-pre-month-weighted average and is invariant to the omitted event-time reference. The last two columns use the stored joint covariance, so they preserve the covariance between the static and dynamic fits. The tests establish nondetection of a difference, not economic equivalence.}
\end{table}


\subsection{Official HonestDiD implementation}

The trend-sensitivity program passes the joint event-time vector and covariance matrix to the official \texttt{HonestDiD} implementation.  The environment receipt fixes HonestDiD 0.2.8 and its source commit, and records the compatible official CVXR and solver versions.  The program does not feed a single static coefficient, $D$, or a standard error into an event-study procedure; its target is $P$.  It reports both the complete 2017Q1--2022Q3 calibration and the prespecified consecutive 2021Q1--2022Q3 post-reopening calibration.  The latter does not delete an internal pandemic quarter or compress a calendar gap.  For the family-conditioned vector its $M=0$ smoothness interval excludes zero, but the interval includes zero at $M=0.005$; the pooled intervals include zero even at $M=0$.  The conventional event-functional interval already includes zero in both structures.  The correct positive zero-exclusion breakdown is therefore undefined: there is no positive robustness threshold to lose.  Earlier installation and API failures remain in the execution log; the successful official run is not replaced with a home-built approximation.

\begin{table}[!htbp]
\centering
\caption{Official HonestDiD sensitivity intervals}\label{tab:app_honestdid}
\small
\begin{tabular}{lllrr}
\toprule
Structure & Calibration & Restriction & Lower & Upper \\
\midrule
Pooled & 2017Q1--2022Q3 & Conventional & $-0.264$ & $0.024$ \\
       &                 & Smoothness, $M=0$ & $-0.140$ & $0.016$ \\
       &                 & Smoothness, $M=0.005$ & $-0.482$ & $0.516$ \\
       &                 & Relative magnitude, $\bar M=0$ & $-0.260$ & $0.019$ \\
       &                 & Relative magnitude, $\bar M=0.5$ & $-1.023$ & $0.789$ \\
Pooled & 2021Q1--2022Q3 & Conventional & $-0.264$ & $0.024$ \\
       &                 & Smoothness, $M=0$ & $-0.230$ & $0.204$ \\
       &                 & Smoothness, $M=0.005$ & $-0.540$ & $0.495$ \\
       &                 & Relative magnitude, $\bar M=0$ & $-0.260$ & $0.019$ \\
       &                 & Relative magnitude, $\bar M=0.5$ & $-0.985$ & $0.742$ \\
\addlinespace
SOC2 monthly paths & 2017Q1--2022Q3 & Conventional & $-0.425$ & $0.010$ \\
       &                 & Smoothness, $M=0$ & $-0.351$ & $0.007$ \\
       &                 & Smoothness, $M=0.005$ & $-0.869$ & $0.352$ \\
       &                 & Relative magnitude, $\bar M=0$ & $-0.420$ & $0.007$ \\
       &                 & Relative magnitude, $\bar M=0.5$ & $-2.222$ & $1.884$ \\
SOC2 monthly paths & 2021Q1--2022Q3 & Conventional & $-0.425$ & $0.010$ \\
       &                 & Smoothness, $M=0$ & $-0.772$ & $-0.007$ \\
       &                 & Smoothness, $M=0.005$ & $-1.100$ & $0.207$ \\
       &                 & Relative magnitude, $\bar M=0$ & $-0.420$ & $0.007$ \\
       &                 & Relative magnitude, $\bar M=0.5$ & $-2.222$ & $1.884$ \\
\bottomrule
\end{tabular}
\tabnote{The target is the observed-month-weighted average of the 15 post-quarter Q5--Q1 coefficients relative to 2022Q4.  Intervals use the full occupation-cluster covariance and official HonestDiD 0.2.8.  Smoothness $M$ is in log points per quarter squared; $\bar M$ bounds post changes relative to the largest signed preperiod change.  Both the full 2017Q1--2022Q3 calibration and the consecutive post-reopening 2021Q1--2022Q3 calibration are fixed and reported.  The recent-window family-path interval excludes zero only at the exact-linear-bias restriction $M=0$ and includes zero at $M=0.005$.  Because each conventional interval includes zero, positive zero-exclusion robustness is undefined, not an omitted threshold.}
\end{table}


\subsection{Onset, endpoint, seasonality, and survey dates}

Onset sensitivity reports every monthly split from November 2022 through June 2023.  January 2023 is retained because December 2022 was designated as the transition month after the November 30, 2022 public release of ChatGPT; robustness across nearby dates does not identify that event as the cause.  The endpoint grid includes December 2024, September and December 2025, and July 2026 wherever the calendar permits.  October 2025 is never imputed.  Relative to the full window, the pooled movements through December 2024, September 2025, and December 2025 are 0.0211, 0.0225, and 0.0212 log point.  Their magnitudes are essentially the same even though their interval endpoints differ.  The full available window remains the canonical target, while December 2024 is reported prominently as the pre-2025-production-break benchmark; the latter is not promoted after outcomes are observed because it retains a negative detected coefficient.  The analogous conditioned movements remain near zero.

\begin{table}[!htbp]
\centering
\caption{Post-outcome exploratory endpoint sensitivity}\label{tab:app_endpoints}
\scriptsize
\setlength{\tabcolsep}{3pt}
\resizebox{\textwidth}{!}{%
\begin{tabular}{lrrrr}
\toprule
Endpoint & Months & Coefficient [95\% CI] & Endpoint $-$ full [paired 95\% CI] & Paired $\mathrm{MDE}_{80}$ \\
\midrule
\multicolumn{5}{l}{\textit{Pooled}} \\
Through December 2024 & 95 & $-0.1110\;[-0.2008,-0.0213]$ & $0.0211\;[-0.0073,0.0495]$ & 0.0407 \\
Through September 2025 & 104 & $-0.1096\;[-0.1978,-0.0214]$ & $0.0225\;[0.0053,0.0397]$ & 0.0245 \\
Through December 2025 (actual gap) & 106 & $-0.1109\;[-0.1979,-0.0239]$ & $0.0212\;[0.0057,0.0367]$ & 0.0221 \\
Through July 2026, excluding September and November 2025 & 111 & $-0.1337\;[-0.2236,-0.0438]$ & $-0.0016\;[-0.0075,0.0043]$ & 0.0084 \\
Through July 2026 (reference) & 113 & $-0.1321\;[-0.2206,-0.0437]$ & Reference & --- \\
\addlinespace
\multicolumn{5}{l}{\textit{SOC2 by calendar month}} \\
Through December 2024 & 95 & $-0.0120\;[-0.1657,0.1417]$ & $0.0097\;[-0.0526,0.0719]$ & 0.0897 \\
Through September 2025 & 104 & $-0.0219\;[-0.1616,0.1178]$ & $-0.0002\;[-0.0345,0.0341]$ & 0.0491 \\
Through December 2025 (actual gap) & 106 & $-0.0170\;[-0.1555,0.1215]$ & $0.0047\;[-0.0248,0.0342]$ & 0.0422 \\
Through July 2026, excluding September and November 2025 & 111 & $-0.0252\;[-0.1660,0.1155]$ & $-0.0035\;[-0.0131,0.0060]$ & 0.0137 \\
Through July 2026 (reference) & 113 & $-0.0217\;[-0.1607,0.1173]$ & Reference & --- \\
\bottomrule
\end{tabular}}
\tabnote{All rows use the rebuilt corrected-preperiod treatment and the same occupation support.  The December 2024 row is the requested pre-2025 endpoint; the December 2025 row retains the actual missing October survey month.  The shutdown-month row excludes September and November 2025 and does not interpolate October.  Nonreference endpoint intervals and paired differences use the common endpoint-grid draws.  To avoid printing a second seed-dependent interval for an identical target, each full-window reference uses its declared canonical single-target interval.  The MDE is the two-sided five-percent, 80-percent normal-theory precision diagnostic, not an economic-equivalence threshold.}
\end{table}


Exposure-group-by-month-of-year interactions and the feasible occupation-specific month-of-year specification both move the target little: the largest absolute movement from the corresponding full-window coefficient is $0.0011$.  A coding-stable post-2020 specification gives $-0.1181$ pooled and $-0.0304$ with family-month effects; paired movements from the full-window estimates are $0.0140$ ($[-0.0288,0.0569]$) and $-0.0087$ ($[-0.1053,0.0878]$).  These replacement fits pass the same numerical certification as the baseline.  The earlier failed attempts and diagnostics remain retained in the failure registry.  Excluding September and November 2025 around the collection gap does not detectably change the full-window estimate.  These rows address sensitivity to calendar construction and survey production, not counterfactual reweighting for the January 2025 or January 2026 population controls.

\subsection{Official-weight and respondent-equivalent era comparison}

Table~\ref{tab:appendix_population_eras} reports a post-outcome exploratory joint-model comparison of 2023--2024 with 2025--2026.  The official-weight row retains the paper's employment-stock estimand; the respondent-equivalent row is a deliberately different routed-count outcome used only to diagnose whether the later pattern mechanically appears only in final-weight stocks.  It cannot recover unpublished age-by-occupation counterfactual weights.

\begin{table}[!htbp]
\centering
\caption{Post-outcome exploratory comparison of early and later postperiods}\label{tab:appendix_population_eras}
\small
\setlength{\tabcolsep}{4pt}
\resizebox{\textwidth}{!}{%
\begin{tabular}{lrrrr}
\toprule
Cell construction & 2023--2024 coefficient [95\% CI] & 2025--2026 coefficient [95\% CI] & Later minus early [95\% CI] & Paired $\mathrm{MDE}_{80}$ \\
\midrule
Official final-weight stocks & $-0.111\;[-0.197,-0.024]$ & $-0.166\;[-0.270,-0.062]$ & $-0.056\;[-0.123,0.012]$ & 0.097 \\
Respondent-equivalent routed counts & $-0.103\;[-0.193,-0.012]$ & $-0.183\;[-0.277,-0.089]$ & $-0.080\;[-0.139,-0.021]$ & 0.085 \\
\bottomrule
\end{tabular}}
\tabnote{This timing comparison was produced after outcomes were opened and is post-outcome exploratory.  Both rows use the corrected 113-month calendar, the same 468 occupations, a joint two-era specification, and 9,999 common occupation multipliers within each paired later-minus-early contrast.  Official final-weight stocks are the paper's employment-stock object.  Respondent-equivalent values are fractional routed-record counts after occupational crosswalking; they are neither distinct survey respondents nor a recovered counterfactual CPS weight vintage.  The official-weight interval does not detect an era difference, whereas the respondent-equivalent interval excludes zero.  That divergence does not identify response composition or show that population-control revisions caused the later estimate.}
\end{table}

